\titleframe

% Contoh Kode Slide
%\section{Introduction}
%\begin{frame}{\insertsectionhead}
%	\framesubtitle{A short introduction to Trigon}
%	\themename is a modern, elegant and versatile theme for Beamer, inspired by
%	the
%	\href{https://github.com/matze/mtheme}{\textsc{metropolis} theme} from Matthias
%	Vogelgesang.
%	\vfill
%	\themename comes with lots of nice extra features
%	\begin{itemize}
%		\item Multiple style variations for title, section and normal slides
%		\item Simple customization of theme colors
%		\item Lots of convenient options to tweak the design
%	\end{itemize}
%\end{frame}

% Contoh Kode Tabel
%  \begin{table}[H]
% \centering
% \caption{A nice table example}
% \begin{tabular}{@{} lccc @{}}
%    \toprule
%    & \textbf{Velocity} & \textbf{Angle}  & \textbf{Vertical force} \\
%    & $U$ & $\alpha$  & $F_z$ \\
%    & [m/s] & [$^\circ$]  & [N] \\
%    \midrule
%    2D simulation  & 9 & 2 & 9.23 \\
%    3D simulation  & 10.0 & 3 & 15.039 \\
%    Experiment A   & 11.31 & 2.5 & 13.2 \\
%    Experiment B   & 11.26 & 2.7 & 12.6 \\
%    Experiment C   & 11.33 & 2.47 & 13.6 \\
%    \bottomrule
%  \end{tabular}
%\end{table}

% Contoh Kode Gambar
%  \begin{figure}[ht!]
%	\begin{subfigure}[b]{0.3\textwidth}
%		\frame{\includegraphics[width=\textwidth]{layout_example-03.jpg}}
%		\caption*{plain}
%	\end{subfigure}
%	\hspace{\fill}
%	\begin{subfigure}[b]{0.3\textwidth}
%		\frame{\includegraphics[width=\textwidth]{layout_example-02.jpg}}
%		\caption*{style1}
%	\end{subfigure}
%	\hspace{\fill}
%	\begin{subfigure}[b]{0.3\textwidth}
%		\frame{\includegraphics[width=\textwidth]{layout_example-01.jpg}}
%		\caption*{style2 (default)}
%	\end{subfigure}
%\end{figure}

% Contoh Kode Blocks
% \begin{block}{Regular block}
%	Just a regular block
%\end{block}
%\begin{alertblock}{Alert block}
%	Some important thing
%\end{alertblock}
%\begin{exampleblock}{Example block}
%	No difference with regular block to avoid excessive distraction
%\end{exampleblock}

\setbeamertemplate{frame footer}{Alauddin Maulana Hirzan, M. Kom - Universitas Semarang}

\section{Teknologi \textit{Smart Home} Berbasis \textit{Internet of Things}}
\begin{frame}{\insertsectionhead}
	\framesubtitle{Apa itu \textit{Smart Home}? \#1}
	\justifying
	Rumah pintar adalah rumah yang dilengkapi dengan teknologi untuk mengontrol dan mengotomatisasi sistem rumah tangga dari jarak jauh seperti penerangan, pintu, termostat, sistem hiburan, alarm keamanan, kamera pengintai, dan peralatan terhubung lainnya.
	\vfill
	Rumah pintar melampaui kendali jarak jauh dan pengaturan yang dapat diprogram yang telah menjadi fitur standar rumah selama beberapa dekade terakhir, untuk menciptakan pemantauan rumah terpadu, kontrol, konservasi energi, ekosistem.
\end{frame}

\begin{frame}{\insertsectionhead}
	\framesubtitle{Apa itu \textit{Smart Home}? \#2}
	\justifying
	Selain fitur, istilah rumah pintar menyiratkan kecerdasan. Rumah Cerdas harus memiliki penginderaan, pemrosesan data, dan fungsi komunikasi nirkabel yang sekunder untuk melaksanakan tugas mereka. 
	\vfill
	Dengan semua fitur tersebut, rumah pintar menjadi mampu mempelajari kebiasaan dan rutinitas penghuninya dan secara otomatis menyesuaikan sistem kontrol rumah dengan cara yang merampingkan tugas-tugas biasa yang menyertai pengoperasian rumah tangga modern. Rumah pintar dibangun dengan konsep membuat kehidupan rumah Anda lebih baik.
\end{frame}

\begin{frame}{\insertsectionhead}
	\framesubtitle{Apa itu \textit{Smart Home}? \#2}
	\justifying
	  \begin{figure}[ht!]
			\begin{subfigure}[b]{0.73\textwidth}
					\frame{\includegraphics[width=\textwidth]{Home-Automation.jpg}}
				\end{subfigure}
		\end{figure}
\end{frame}

\begin{frame}{\insertsectionhead}
	\framesubtitle{Apa itu \textit{Smart Home}? \#3}
	\justifying
	Dengan menggunakan teknologi \textit{Internet of Things}, pemilik rumah dapat merasakan  manfaat berupa:
	\vfill
	\begin{itemize}
		\item Hemat energi dengan termostat pintar
		\item Kontrol akses pintu
		\item Pantau rumah melalui video
		\item Atur suasana hati dengan lampu pintar
		\item Kontrol pintu garasi dari mana saja
		\item Kelola rumah dari hub pintar
		\item Kendali suara
		\item Pastikan Rumah Tetap Aman
	\end{itemize}
\end{frame}

\begin{frame}{\insertsectionhead}
	\framesubtitle{Membangun \textit{Smart Home}? \#1}
	\justifying
	Untuk membangun \textit{Smart Home} yang aman dan nyaman, tentunya tidak bisa dilakukan semudah itu. Ada beberapa tahapan yang harus dipenuhi berupa:
	\vfill
	\begin{itemize}
		\item Memilih Penyedia Layanan (baik Internet maupun sistem otomatisasi)
		\item Mendesain \textit{Smart Home} sesuai dengan teknologi yang ingin dipasang
		\item Pemasangan Perangkat
		\item Pengujian Sistem
	\end{itemize}
\end{frame}

\begin{frame}{\insertsectionhead}
	\framesubtitle{Membangun \textit{Smart Home}? \#2 - Penyedia Layanan}
	\justifying
	Aspek ini merupakan aspek paling utama yang tidak bisa ditinggalkan. Karena tanpa adanya penyedia layanan, perangkat-preangkat IoT tidak akan bisa melakukan pengiriman data maupun otomatisasi. DI aspek ini yang merupakan penyedia layanan yang harus diperhatikan adalah
	\vfill
	\begin{itemize}
		\item Layanan Internet
		\item Layanan Otomatisasi
	\end{itemize}
\end{frame}

\begin{frame}{\insertsectionhead}
	\framesubtitle{Membangun \textit{Smart Home}? \#3 - Layanan Internet \#1}
	\justifying
	Layanan ini merupakan hal vital yang menjadi syarat utama untuk membuat \textit{Smart Home}. Tanpa ada nya layanan ini perangkat tidak akan bisa berkomunikasi satu sama lainnya. Sehingga pengendalian maupun pemantauan menjadi mustahil untuk dilakukan.
	\vfill
	Pemasangan layanan internet tidak harus menggunakan layanan kencang, namun sesuai dengan kebutuhan perangkat-perangkat yang akan dipakai. Karena sebagian besar perangkat \textbf{belum tentu} memakan banyak \textit{bandwidth}
\end{frame}

\begin{frame}{\insertsectionhead}
	\framesubtitle{Membangun \textit{Smart Home}? \#4 - Layanan Internet \#2}
	\justifying
	Contoh:
	\vfill
	\begin{quote}
		Penggunaan MQTT maupun protokol lain lebih hemat dibandingkan pengguna membuka media sosial apalagi media streaming seperti Youtube
	\end{quote}
	\vfill
	\begin{block}{Info}
		Jika satu orang melakukan streaming video sebesar 100MB, maka ukuran tersebut cukup untuk mengirimkan  149.369 kali perintah (termasuk konektivitas TCP dan data pemantauan).  Sumber: https://docs.devicewise.com/Content/Products/IoT\_Portal\_API\_Reference\_Guide/MQTT
		\_Interface/MQTT-data-usage.htm
	\end{block}
\end{frame}

\begin{frame}{\insertsectionhead}
	\framesubtitle{Membangun \textit{Smart Home}? \#5 - Layanan Otomatisasi \#1}
	\justifying
	Aspek ini bersifat operasional dan tergantung dari sistem yang dibuat. Beberapa pengguna lebih suka membuat otomatisasi mereka sendiri. Namun juga ada yang lebih mengandalkan layanan yang sudah ada seperti:
	\vfill
	\begin{itemize}
		\item IFTTT
		\item Google Nest
		\item Blynk
		\item Cayenne
		\item KaaIoT
		\item dll
	\end{itemize}
\end{frame}

\begin{frame}{\insertsectionhead}
	\framesubtitle{Membangun \textit{Smart Home}? \#6 - Layanan Otomatisasi \#2 - IFTTT}
	\justifying
	Dikutip dari \textbf{ifttt.com} IFTTT adalah
	\vfill
	\begin{quote}
		Kependekan dari \textbf{If This Then That}, dan merupakan cara terbaik untuk mengintegrasikan aplikasi, perangkat, dan layanan. Teknologi ini membantu perangkat, layanan, dan aplikasi bekerja sama dengan cara baru dan canggih.
		\vfill
		IFTTT didirikan pada keyakinan bahwa segala sesuatu bekerja lebih baik bersama-sama. Ketidakcocokan teknologi telah menjadi tantangan bagi siapa pun yang mencoba membangun rumah pintar atau membuat rutinitas otomatis dalam hidup mereka. IFTTT membuatnya mudah.
	\end{quote}
\end{frame}

\begin{frame}{\insertsectionhead}
	\framesubtitle{Membangun \textit{Smart Home}? \#7 - Layanan Otomatisasi \#3 - IFTTT}
	\justifying
	Apa yang bisa IFTTT lakukan?
	\vfill
	Sederhananya, IFTTT bisa melakukan apa saja! Beberapa Applet rumah pintar yang paling populer termasuk menggunakan cuaca untuk menentukan suhu di rumah, menyalakan sistem keamanan saat pengguna meninggalkan rumah, dan menggunakan suara untuk menya-lakan perangkat.
	\vfill
	\begin{figure}[ht!]
		\begin{subfigure}[b]{0.4\textwidth}
			\frame{\includegraphics[width=\textwidth]{ifttt.png}}
		\end{subfigure}
	\end{figure}
\end{frame}

\begin{frame}{\insertsectionhead}
	\framesubtitle{Membangun \textit{Smart Home}? \#8 - Layanan Otomatisasi \#4 - IFTTT}
	\justifying
	\begin{figure}[ht!]
		\begin{subfigure}[b]{0.7\textwidth}
			\frame{\includegraphics[width=\textwidth]{ifttt2.png}}
		\end{subfigure}
	\end{figure}
\end{frame}

\begin{frame}{\insertsectionhead}
	\framesubtitle{Membangun \textit{Smart Home}? \#9 - Layanan Otomatisasi \#5 - Layanan Lain}
	\justifying
	Layanan lain juga menyediakan layanan otomatisasi mereka masing-masing, selain itu mereka juga menyediakan fitur penyimpanan data pemantauan sementara dengan gratis untuk beberapa perangkat.
	\vfill
	\begin{figure}[ht!]
		\begin{subfigure}[b]{0.43\textwidth}
			\frame{\includegraphics[width=\textwidth]{blynk.png}}
		\end{subfigure}
	\end{figure}
\end{frame}

\begin{frame}{\insertsectionhead}
	\framesubtitle{Membangun \textit{Smart Home}? \#10 - Desain dan Pasang Perangkat}
	\justifying
	Di tahapan ini memerlukan desain dasar rumah agar membantu proses pemasangan perangkat. Letak dinding, furniture, maupun benda-benda lainnya mampu mempengaruhi pemasangan perangkat.
	\vfill
	\begin{figure}[ht!]
		\begin{subfigure}[b]{0.5\textwidth}
			\frame{\includegraphics[width=\textwidth]{layout.jpg}}
		\end{subfigure}
	\end{figure}
\end{frame}

\begin{frame}{\insertsectionhead}
	\framesubtitle{Membangun \textit{Smart Home}? \#11 - Pengujian}
	\justifying
	Tahapan ini memasikan semua perangkat bekerja dengan baik. Baik perangkat pemantauan hingga pengamanan rumah. Pengujian dapat dilakukan dengan metode standar seperti:
	\vfill
	\begin{itemize}
		\item \textit{Black Box}
		\item \textit{White Box}
		\item \textit{Internet of Things Test Suites}
	\end{itemize}
\end{frame}

% Bagian Kedua

\section{Teknologi \textit{Smart City} Berbasis \textit{Internet of Things}}
\begin{frame}{\insertsectionhead}
	\framesubtitle{Mengenal \textit{Smart City}}
	\justifying
	Berdasarkan \textbf{cisco.com}, \textit{Smart City} adalah
	\vfill
	\begin{quote}
		Sebuah kota pintar menggunakan teknologi digital untuk menghubungkan, melindungi, dan meningkatkan kehidupan warga. Sensor IoT, kamera video, media sosial, dan input lainnya bertindak sebagai sistem saraf, memberikan umpan balik yang konstan kepada operator kota dan warga sehingga mereka dapat membuat keputusan yang tepat.
	\end{quote}
	\vfill
	Sumber : https://www.cisco.com/c/en/us/solutions/industries/smart-connected-com\\
	munities/what-is-a-smart-city.html
\end{frame}

\begin{frame}{\insertsectionhead}
	\framesubtitle{Karakteristik \textit{Smart City} \#1}
	\justifying
	Apa faktor-faktor yang menjadikan sebuah kota sebuah \textit{Smart City}? Perhatikan parameter berikut:
	\vfill
	\begin{itemize}
		\item \textit{Smart Economy}
		\item \textit{Smart People}
		\item \textit{Smart Governance}
		\item \textit{Smart Mobility}
		\item \textit{Smart Environment}
		\item \textit{Smart Living}
	\end{itemize}
\end{frame}

\begin{frame}{\insertsectionhead}
	\framesubtitle{Karakteristik \textit{Smart City} \#2 - \textit{Smart Economy}}
	\justifying
	Ekonomi sering mengacu pada pengelolaan kekayaan dan sumber daya kota, terutama dalam hal produksi dan konsumsi barang dan jasa. Faktor-faktor yang harus dimasukkan dalam ekonomi cerdas adalah kewirausahaan, inovasi, merek dagang, fleksibilitas pasar tenaga kerja, produktivitas, integrasi di pasar internasiona, dan lain-lain.
	\vfill
	\begin{block}{Info}
		Hal ini secara otomatis mempengaruhi kesejahteraan warga yang tinggal di kota tersebut
	\end{block}	
\end{frame}

\begin{frame}{\insertsectionhead}
	\framesubtitle{Karakteristik \textit{Smart City} \#3 - \textit{Smart People}}
	\justifying
	Kota Pintar akan terus tumbuh dan matang selama ada orang pintar dan teknologi pintar yang mendukungnya. Orang pintar dapat diukur dari beberapa indikator seperti pendidikan, kreativitas, inovasi, partisipasi, dan lain-lain
	\vfill 
	Tingkat pendidikan penduduk dapat dengan mudah diukur dari tingkat pendidikan yang ditempuh. Pendidikan menengah, perguruan tinggi, atau gelar pendidikan lain dalam suatu populasi.
	\vfill
	\begin{block}{Info}
		IQ, EQ, dan SQ bisa menjadi faktor dalam menentukan Orang Pintar dalam suatu populasi
	\end{block}
\end{frame}

\begin{frame}{\insertsectionhead}
	\framesubtitle{Karakteristik \textit{Smart City} \#4 - \textit{Smart Governance}}
	\justifying
	Aspek ini merujuk ke dalam organisasi yang mengatur suatu populasi (pemerintahan). Tata Kelola yang dilakukan oleh pemerintahan dapat diartikan sebagai cara untuk sebuah kota diatur/diorganisasikan secara internal. Setiap \textit{Smart City} memiliki cara yang berbeda-beda tergantung dari pemerintahan yang mengatur. Sehingga dapat mempengaruhi dinamika ekonomi.
	\vfill
	Pemerintahan yang baik memiliki tiga indikator:
	\vfill
	\begin{itemize}
		\item Keterbukaan/Transparansi
		\item Kolaborasi Efektif
		\item Berbasis Kinerja (\textit{Meritocracy})
	\end{itemize}
\end{frame}

\begin{frame}{\insertsectionhead}
	\framesubtitle{Karakteristik \textit{Smart City} \#5 - \textit{Smart Mobility} \#1}
	\justifying
	Aspek ini merujuk ke bagaimana warga di sebuah kota berpindah dari satu tempat ke tempat lain. Transportasi juga menjadi indikator dari \textit{Smart City} karena mempengaruhi teknologi yang digunakan.
	\vfill
	Moda transportasi yang dimaksud terintegrasi sehingga memudahkan warga untuk menggunakan moda-moda transportasi yang disediakan kota.
	\vfill
	\begin{block}{Info}
		\textbf{\textit{Terminal Bersepadu Selatan}} (TBS) di Kuala Lumpur merupakan contoh nyata integrasi moda transportasi Kereta Api, Bus, Taxi, dan lain-lain
	\end{block}
\end{frame}

\begin{frame}{\insertsectionhead}
	\framesubtitle{Karakteristik \textit{Smart City} \#6 - \textit{Smart Mobility} \#2}
	\justifying
	\begin{figure}[ht!]
		\begin{subfigure}[b]{0.65\textwidth}
			\frame{\includegraphics[width=\textwidth]{TBS.jpg}}
		\end{subfigure}
	\end{figure}
\end{frame}

\begin{frame}{\insertsectionhead}
	\framesubtitle{Karakteristik \textit{Smart City} \#7 - \textit{Smart Environment}}
	\justifying
	\textit{Green Communication} atau disebut dengan \textbf{Komunikasi Hijau} adalah salah satu tema umum yang diterapkan dalam Kota Pintar. Kota-kota ini ingin mengurangi polusi karbon yang ada di lingkungan mereka.
	\vfill
	Terdapat beberapa upaya yang dilakukan dengan cara mengganti ke kendaraan yang lebih ramah lingkungan (Kendaraan Listrik), Pengelolaan Sampah yang efisien (Pengolahan Ulang Limbah) hingga bangunan pintar yang dilengkapi dengan panel surya untuk menekan penggunaan listrik berbasis fosil.
\end{frame}

\begin{frame}{\insertsectionhead}
	\framesubtitle{Karakteristik \textit{Smart City} \#8 - \textit{Smart Living}}
	\justifying
	Penduduk merupakan kunci kemajuan suatu kota. Peningkatan kualitas hidup penduduk sangat penting dan perlu juga memperhatikan fasilitas kebudayaan, kondisi kesehatan, keselamatan individu, kualitas perumahan, fasilitas pendidikan, daya tarik wisata, dan kohesi (relasi) sosial.
	\vfill
	\begin{block}{Info}
		Kota yang baik memfasilitasi semua warganya tanpa pandang bulu baik secara ras, agama, suku, bahkan warna kulit.
	\end{block}
\end{frame}

\begin{frame}{\insertsectionhead}
	\framesubtitle{\textit{What to do?}}
	\justifying
	Implementasi \textit{Smart City} bukanlah hal yang mudah. Selain faktor individu warga, peran pemerintahan yang tepat dapat membantu proses pembangunan \textit{Smart City}. Jika faktor-faktor tersebut tidak terpenuhi, maka \textit{Smart City} hanyalah mimpi belaka yang tidak dapat memberikan keadilan yang tepat ke semua warganya.
	\vfill
	\begin{quote}
		\centering
		Smart City Are Built By Smart People, Not Smart Things\\
		\raggedleft -forbes.com
	\end{quote} 
\end{frame}

\setbeamertemplate{background} 
{
	\includegraphics[width=\paperwidth,height=\paperheight]{thank-you.jpg}
}
\begin{frame}[plain]
\end{frame}